\chapter{A molecular biology primer}\label{app:biology}
\markboth{A molecular biology primer}{A molecular biology primer}

This appendix gives a reader from computer science, mathematics or physics the
biological vocabulary the book assumes. It is a map, not a substitute for the
chapters: each section points to the chapter where the mechanism is developed
properly. Figure~\ref{fig:primer-helix} in Chapter~1 shows the structures
referred to here.

\section{Nucleotides, strands and the double helix}
DNA (deoxyribonucleic acid) is a polymer. Its repeating unit, the
\Term{nucleotide}, has three parts: a five-carbon sugar (deoxyribose), a
phosphate group and one of four \Term{bases}: adenine (A), cytosine (C),
guanine (G) or thymine (T). Phosphates link the sugar of one nucleotide to the
sugar of the next, forming a covalent \emph{sugar--phosphate backbone}. Because
the link always joins the $3'$ carbon of one sugar to the $5'$ carbon of the
next, a strand has two chemically different ends, called the $5'$ end and the
$3'$ end, and therefore a direction. By convention sequences are written from
$5'$ to $3'$.

Two strands with complementary sequences can pair: A with T and G with C. The
paired strands run in opposite directions (they are \Term{antiparallel}) and
wind into the right-handed double helix described by Watson and Crick in
1953.\cite{WatsonCrick1953} A base pair is held by hydrogen bonds (two for
A$\cdot$T, three for G$\cdot$C), but much of the duplex's stability comes from
\emph{stacking} between neighbouring base pairs. Pairing is reversible: heating
separates the strands (\emph{melting}), and cooling lets complementary strands
find each other again (\emph{annealing} or \emph{hybridisation}).

Three distinctions matter throughout the book.
\begin{itemize}
\item A \emph{base} is a chemical group; a \emph{nucleotide} is base plus sugar
plus phosphate; a \emph{strand} is a covalent chain of nucleotides; a
\emph{duplex} is two strands held together by pairing.
\item The \emph{complement} of a sequence replaces each base by its partner; the
\emph{reverse complement} also reverses the order, so that the partner strand is
again written $5'\to3'$. The partner of $5'$-ACGTTG-$3'$ is
$5'$-CAACGT-$3'$ (Chapter~5).
\item Pairing is not a covalent bond. Two strands lying end to end on a
template are still two molecules until an enzyme (a ligase) joins them.
\end{itemize}

\section{DNA, RNA and proteins}
Cells use three classes of information-carrying polymer
(Figure~\ref{fig:central-dogma}). DNA stores genetic information. RNA
(ribonucleic acid) differs from DNA by a hydroxyl group on its sugar and by
using uracil (U) in place of thymine; it is usually single-stranded and often
folds back on itself. Proteins are chains of amino acids; they fold into
three-dimensional machines, including \Term{enzymes}, which catalyse chemical
reactions.

\TikzPlate{backmatter/figures/central-dogma}{Information flow in a cell and
the molecular machines that carry it out. RNA polymerase (gold) reads a DNA
template strand and synthesises an RNA copy in the $5'\to3'$ direction
(transcription); a ribosome reads the RNA three bases at a time and joins
amino acids into a protein (translation). A transcription factor binding near
a gene can raise or lower how often that gene is transcribed, which is the
basis of the regulatory circuits in Chapter~30. Molecules are drawn
schematically and not to scale.}{fig:central-dogma}

The flow DNA $\to$ RNA $\to$ protein is called the \Term{central dogma}.
\emph{Transcription} copies a stretch of DNA (a gene) into RNA; \emph{translation}
reads the RNA in three-letter codons to build a protein. Proteins called
\Term{transcription factors} bind DNA near genes and switch transcription up or
down. Gene regulation is therefore a natural information-processing system:
inputs (signals, factor concentrations) change outputs (which proteins are
made), and feedback can create memory. Chapter~30 develops this view
quantitatively, and Chapter~32 asks where the computational analogy holds.

\section{Enzymes that act on DNA}
Most laboratory DNA manipulations use a small set of enzymes, each of which
recognises a particular substrate and changes one specific chemical feature
(Chapter~8).
\begin{description}
\item[Polymerases] extend a primer that is paired to a template, adding
nucleotides only to the $3'$ end, so new DNA grows $5'\to3'$. Heat-stable
polymerases make PCR possible (Chapter~9).
\item[Ligases] seal a \emph{nick}, joining a $3'$ hydroxyl to an adjacent $5'$
phosphate on the same template, so that two strands become one.
\item[Restriction enzymes] (restriction endonucleases) recognise a short
sequence such as GAATTC and cut both strands at defined positions, often leaving
short single-stranded overhangs (``sticky ends'').
\item[Exonucleases] remove nucleotides from strand ends; \emph{endonucleases}
cut within a strand. Nucleases in serum and cells degrade unprotected DNA
circuits (Chapter~31).
\item[CRISPR-associated nucleases] such as Cas9 are guided by an RNA to a
complementary DNA site and cut there; engineered variants write information into
genomes and are used as molecular recorders (Chapter~32 and the Epilogue).
\end{description}
An enzyme is not a perfect logical operator. It must find its substrate,
bind, react and release; it works at finite rate, sometimes acts on near-miss
sequences, and depends on temperature, salt and cofactors such as ATP or
Mg$^{2+}$.

\section{The laboratory toolkit at a glance}
Figure~\ref{fig:lab-toolkit} summarises the operations that recur throughout
the book.

\TikzPlate{backmatter/figures/lab-toolkit}{The molecular programmer's laboratory
toolkit. \emph{Write}: short strands (oligonucleotides, typically 20--200
nucleotides) are made by chemical synthesis. \emph{Mix and react}: strands are
combined in buffer at controlled temperature, where they pair, displace one
another or are acted on by enzymes. \emph{Copy}: PCR amplifies a chosen interval
exponentially. \emph{Separate}: gel electrophoresis sorts molecules mainly by
length; beads with capture probes pull out molecules carrying a sequence.
\emph{Read}: fluorescence reports bulk changes in real time; sequencing reads
individual molecules' sequences. Each arrow is an operation with a yield, an
error rate and a cost, not an ideal function.}{fig:lab-toolkit}

\begin{description}
\item[Chemical synthesis] builds a strand one nucleotide at a time on a solid
support. Each coupling step succeeds with high but imperfect probability, so
long strands accumulate errors and truncated products; this is why designs
favour many short strands. Enzymatic and array-based synthesis are lowering the
cost of writing large libraries.\cite{Kosuri2014}
\item[Annealing] is a heat-then-cool protocol that lets strands find their
lowest-free-energy partners (Chapter~6).
\item[PCR] (polymerase chain reaction) repeats denaturation, primer annealing and
extension; each cycle can double the number of copies of the interval between
two primers (Chapter~9).\cite{Saiki1988}
\item[Gel electrophoresis] pulls negatively charged DNA through a gel; shorter
fragments move further. A band reports a size class, not a sequence
(Chapter~10).
\item[Fluorescence] reports molecular events optically. A common trick in DNA
circuits places a fluorophore and a quencher close together on a duplex, so
that releasing an output strand separates them and the solution brightens.
\item[Sequencing] reads the base order of individual molecules. Massively
parallel short-read sequencing reads millions of fragments; nanopore sequencing
reads long single molecules by the current changes they cause in a pore.
Sequencing is the readout of DNA data storage (Chapter~25).
\end{description}

\section{Numbers worth remembering}
Molecular computing is quantitative, and a feeling for magnitudes prevents
many errors.
\begin{center}
\small
\begin{tabular}{@{}p{0.43\linewidth}p{0.5\linewidth}@{}}
\toprule
Quantity & Typical value\\
\midrule
Rise per base pair in B-DNA & 0.34 nm\\
Base pairs per helical turn & about 10.5\\
Duplex diameter & about 2 nm\\
Mass of one nucleotide in a strand & about 330 g\,mol$^{-1}$ (Chapter~4)\\
Molecules in 1\,\textmu L of a 1\,\textmu M solution & $6.0\times10^{11}$\\
Typical circuit concentrations & 1\,nM--1\,\textmu M\\
Hybridisation rate constant (good toehold) & about $10^{6}$\,M$^{-1}$s$^{-1}$ (Chapter~19)\\
Free energy scale $RT$ at 25\,\textdegree C & 2.48\,kJ\,mol$^{-1}$ = 0.593\,kcal\,mol$^{-1}$\\
\bottomrule
\end{tabular}
\end{center}
The molecule count follows from $N = c\,V\,N_{\mathrm A}$:
$10^{-6}\,\mathrm{mol\,L^{-1}}\times10^{-6}\,\mathrm{L}\times
6.022\times10^{23}\,\mathrm{mol^{-1}}=6.0\times10^{11}$. Even a tiny volume holds
an enormous number of molecules, which is the source of molecular parallelism;
Chapter~4 shows why that number is nevertheless finite and quickly exhausted by
exponential search.

\section{Further reading}
Any modern molecular biology textbook covers this material in depth; a
freely accessible chapter on DNA structure is available in \emph{Molecular
Biology of the Cell}.\cite{Alberts2002} The review by Seeman and Sleiman
introduces DNA as an engineering material.\cite{SeemanSleiman2018}
